% Part: lambda-calculus
% Chapter: introduction
% Section: lambda-computable

\documentclass[../../../include/open-logic-section]{subfiles}

\begin{document}

\olfileid{lam}{int}{cmp}
\olsection{Fungsi kang \usetoken{S}{lambda definable} Bisa Diitung}

\begin{thm}
\ollabel{thm:lambda-computable}
Yen fungsi parsial $f$ !!{lambda defined}
dening suku lambda, fungsi iku bisa diitung.
\end{thm}

\begin{proof}
Upamana fungsi~$f$ !!{lambda defined} dening
suku lambda~$X$. Kita andharake prosedur
informal kanggo ngitung~$f$. Kanggo input
$m_0$, \dots,~$m_{n-1}$, tulisen suku
$X \num{m_0} \ldots \num{m_{n-1}}$.
Gawea sawijining wit: luwih dhisik tulisen
kabeh asil reduksi saklangkah saka suku wiwitan;
ing sangisore iku tulisen kabeh asil reduksi
saklangkah saka suku-suku mau (yaiku asil
reduksi rong langkah saka suku wiwitan);
banjur terusna. Yen ing sawijining wektu
ketemu numeral, wangsulana wilangan kang
diwakili numeral iku; yen ora nate ketemu,
fungsi iku ora nduweni nilai ing input kasebut.

Kanthi nggunakake tesis Church, kita bisa
nyimpulake yen fungsi iki bisa diitung.
Cara kang luwih becik kanggo mbuktekake
teorema iki yaiku menehi katrangan rekursif
babagan prosedur panelusuran mau. Umpamane,
kita bisa netepake sawetara fungsi lan relasi
rekursif primitif, ``$\fn{IsASubterm}$,''
``$\fn{Substitute}$,''
``$\fn{ReducesToInOneStep}$,''
``$\fn{ReductionSequence}$,''
``$\fn{Numeral}$,'' lan sateruse.
Prosedur rekursif parsial kanggo ngitung
$f(m_0, \dots, m_{n-1})$ banjur nggoleki
urutan reduksi saklangkah kang diwiwiti saka
$X \num{m_0} \dots \num{m_{n-1}}$ lan
dipungkasi dening sawijining numeral,
banjur mbalekake wilangan kang cocog karo
numeral iku. Rinciane dawa lan mboseni,
nanging saliyane iku lumrah.
\end{proof}

\end{document}
